\lecture{10}{2026-10-20}{Elliptic Curves, Algebraically}

While we defined the group operation on an elliptic curve geometrically, we can also define it algebraically.

Given an elliptic curve \(E\), we denote by \(E(K)\) the set of points on \(E\) with coordinates in a field \(K\), together with the point at infinity \(\infty\).

\begin{definition}[Group Operation on an Elliptic Curve, Algebraically]
    Let \(K\) be a field, and let \(E\) be an elliptic curve defined over \(K\) given by the Weierstrass equation
    \[
        E \colon y^2 = x^3 + Ax + B.
    \]
    For any two points \(P\) and \(Q\) on \(E\), we define their sum \(P + Q\) as follows:
    \begin{enumerate}
        \item If \(P = \infty\), then \(P + Q = Q\). If \(Q = \infty\), then \(P + Q = P\).
        \item If \(P = (x_1, y_1)\) and \(Q = (x_2, y_2)\) with \(x_1 = x_2\) and \(y_1 = -y_2\), then \(P + Q = \infty\).
        \item If \(P = (x_1, y_1)\) and \(Q = (x_2, y_2)\) with \(x_1 \neq x_2\) or \(y_1 \neq -y_2\), we define \(P + Q = (x_3, y_3)\) where
        \begin{align*}
            m &=
            \begin{cases}
                \left(y_2 - y_1\right) / \left(x_2 - x_1\right) & \text{if } P \neq Q, \\
                \left(3x_1^2 + A\right) / \left(2y_1\right) & \text{if } P = Q,
            \end{cases} \\
            x_3 &= m^2 - x_1 - x_2, \\
            y_3 &= m(x_1 - x_3) - y_1.
        \end{align*}
    \end{enumerate}
\end{definition}

% somethign something on how to pick (A, B, p)

% theorems by hasse, lenstra, and schoof for sampling.

% another (bad) way is to choose A = 0 and p + 1 = 2p'. sam will tell us why it's bad later.

\begin{definition}[Torsion]
    The \emph{\(n\)-torsion subgroup} of \(E(K)\) is 
    \[
        E(K)[n] = \{ P \in E(K) : nP = \infty \}.
    \]
\end{definition}

\begin{example}
    For \(n = 2\), for any point \(P = (x, y) \in E(K)\), we have \(2P = \infty\) if and only if \(P = -P\). This occurs if and only if \(y = 0\). Thus, the 2-torsion points are the points on the curve where the \(y\)-coordinate is zero, along with the point at infinity.

    How many \(x\) satisfy \(0 = x^3 + Ax + B\)? The polynomial \(x^3 + Ax + B\) has either \(0\), \(1\), or \(3\) roots in \(K\). Therefore, the 2-torsion subgroup \(E(K)[2]\) can have either \(1\), \(2\), or \(4\) points, depending on the number of roots of the polynomial in the field \(K\), and therefore it is isomorphic to either \(\{0\}\), \(\mathbb{Z}_2\), or \(\mathbb{Z}_2 \times \mathbb{Z}_2\).
\end{example}

\begin{theorem}
    Let \(p\) be a prime, and let \(n\) not be divisible by \(p\). Then
    \(E(\overline{\mathbb{F}_p})[n] \cong \mathbb{Z}_n \times \mathbb{Z}_n\).
\end{theorem}

\begin{corollary}
    Let \(p\) be a prime. Then \(E(\mathbb{F}_p) = \mathbb{Z}_{m_1} \times \mathbb{Z}_{m_2}\) for some \(m_1, m_2\).
\end{corollary}

% --- I couldn't keep track

\begin{definition}[Pairing]
    A \emph{pairing} is a non-degenerate bilinear map
    \[
        e \colon E(\mathbb{F}_p)[n] \times E(\mathbb{F}_p)[n] \to \mu_n,
    \]
    where \(\mu_n\) is a subgroup of \(\mathbb{F}_{p^k}\).
\end{definition}

Here, bilinear means that for any \(P, Q, R \in E(\mathbb{F}_p)[n]\) and \(a, b \in \mathbb{Z}\), we have
\[
    e(aP + bQ, R) = e(P, R)^a \cdot e(Q, R)^b \quad \text{and} \quad e(P, aQ + bR) = e(P, Q)^a \cdot e(P, R)^b.
\]

One cool property is that
\[
    e(iP, jQ) = e(P, Q)^{ij} \quad \text{for any } i, j \in \mathbb{Z}.
\]

Three cool things:

Check equality of discrete log:
Given \(P_1, Q_1 = x_1 P_1\) and \(P_2, Q_2 = x_2 P_2\), we can check if \(x_1 = x_2\) by checking if \(e(P_1, Q_2) = e(Q_1, P_2)\), since
\[
    e(P_1, Q_2) = e(P_1, x_2 P_2) = e(P_1, P_2)^{x_2} \text{ and } e(Q_1, P_2) = e(x_1 P_1, P_2) = e(P_1, P_2)^{x_1}.
\]

Second cool thing.
Triple diffie helman.
Parameter is a point \(P\), Alice, Bob, and Charlie each have a private key \(a, b, c\) and a public key \(A = aP, B = bP, C = cP\). They can compute a shared secret \(e(A, B)^c = e(A, C)^b = e(B, C)^a = e(P, P)^{abc}\).

Third cool thing.
Subexponential DLOG attacks.
Given an elliptic curve DLOG problem, \(P, Q = xP\), we can compute \(e(P, P)\) and \(e(P, Q) = e(P, P)^x\). Then we can solve the DLOG problem in \(\mathbb{F}_{p^k}\) using subexponential algorithms like index calculus.

For this attack to not work, we need \(k\) to be large enough.
But it would also be nice if \(k\) was small enough so that we could use the other two interesting properties.